\section{rCFD and CFD (Prof.\ Pirker)}

\begin{summary}{rCFD in two minutes}
\textsf{One idea.} Solve the expensive flow once, record \emph{how the fluid moves} as
cell-to-cell maps, then replay those maps many times to carry a scalar (species, temperature).
Momentum and pressure are never solved again: hundreds of times faster.
\[
T_j^{*}=\sum_iM^{(f)}_{ji}T_i,\qquad
\mathbf{T}^{n+1}=\mathcal{B}\Bigl(\mathcal{C}\bigl(\mathcal{S}(\mathbf{M}^{(f_n)}\mathbf{T}^n)\bigr)\Bigr),\qquad
R_{fg}=\lVert\Phi^{(f)}-\Phi^{(g)}\rVert .
\]
map $\mathbf{M}$ (advection) \textbullet{} face swap $\mathcal{S}$ (diffusion) \textbullet{}
correction $\mathcal{C}$ (your work) \textbullet{} conservation constraint $\mathcal{B}$.\
\textsf{Pirker's three bold failures:} (1) mass imbalance $\to$ global monitor + mass-acting face
swaps; (2) unphysical diffusion $\to$ locally adapted swaps (sub-grid viscosity); (3) weak
recurrence prominence $\to$ recurrence islands.\
\textsf{Works:} recurring flow, passive scalar. \textsf{Fails:} the scalar drives the flow, or the
process changes state $\to$ multi-state rCFD.
\end{summary}

\pic{A GPS that recorded one drive through the city. Next time you only replay the recorded route
(cheap), you do not recompute it. It works as long as the streets are the same. If a street is
closed (a new boundary condition, like the cooled lid), the recording does not know the detour:
you need a new recording for that situation.}

\board{$\mathbf{T}^{n+1}=\mathcal{B}(\mathcal{C}(\mathcal{S}(\mathbf{M}\mathbf{T}^n)))$ \textbullet{}
$\Delta_{\mathrm{swap}}=\fsd V\,(T_{c_1}-T_{c_0})/2$, $D_{\mathrm{eq}}=\fsd c_s^2/(2\Delta t)$ \textbullet{}
limiter $\Delta q_f=s\min(Q_h^+,Q_c^-)$ \textbullet{}
$Co=u\Delta t/\Delta x<1$, $d=\Gamma\Delta t/\Delta x^2<0.5$ \textbullet{} $Re=uL/\nu$.}

\begin{qabox}[Pirker]{\pred Flow-based versus transport-based rCFD?}
\from{Numerical Methods in Fluid Mechanics (Pirker) · rCFD lecture, slides 16--34 · predicted}
\ans Flow-based: stitch recorded flow sequences into a long artificial flow and solve a transport
equation on it (Lichtenegger \& Pirker 2016, about 100$\times$ faster). Transport-based: replace
convection by cell-to-cell shifts and diffusion by face swaps, no equation solved (Pirker \&
Lichtenegger 2018, steel ladle 26~h $\to$ 1.6~s). My thesis uses the transport-based form.
\end{qabox}

\begin{qabox}[Pirker]{\pred Recurrence matrix and prominence -- and your case?}
\from{Numerical Methods in Fluid Mechanics (Pirker) · rCFD slides 50--51; thesis Sec.~2.4 · predicted}
\ans The recurrence norm compares two recorded states; all pairs form the matrix. A row (recurrence
vector) rises and oscillates about a plateau; deep dips = the flow comes back = we can jump.
Prominence $P_{\mathrm{rel}}=\langle d_{\max}-d_{\min}\rangle/\bar d$, required 0.05--0.1. Mine: median
0.028, two blocks (filling, settling), no return $\to$ sequential replay.
\end{qabox}

\begin{qabox}[Pirker]{\pred Why did the recorded local diffusivity not help?}
\from{Thesis Sec.~2.8.1; Numerical Methods in Fluid Mechanics (Pirker) · explicit limits · predicted}
\ans The exchange is explicit and capped at half a cell per face and step, a diffusion number
$\mathrm{Di}=D\Delta t/c_s^2\le0.5$. The recorded $D$ asks for $\mathrm{Di}\approx4$ everywhere, so
every face hits the cap and the field behaves like $\fsd=0.5$ (three replays bit-identical).
An implicit exchange or a smaller step would let it act. This is the explicit stability limit of
the lecture, $d<0.5$.
\end{qabox}

\begin{qabox}[Pirker]{\pred Numerical diffusion: where in your thesis?}
\from{Numerical Methods in Fluid Mechanics (Pirker) · FV~II numerical diffusion; thesis Sec.~4.2 · predicted}
\ans First-order upwind is stable but adds false diffusion that smears fronts. In my replay the
smearing comes from the face exchange, not the maps (maps only keep the front at 2--4~mm). My
correction is an anti-diffusive step bounded like flux-corrected transport: conservative, no new
extremum, front 15.7 $\to$ 5.3/4.5~mm against 5.0~mm.
\end{qabox}

\begin{qabox}[Pirker]{\pred rCFD and digital twins.}
\from{Numerical Methods in Fluid Mechanics (Pirker) · rCFD slides 59--72 digital twins · predicted}
\ans CSTR models are instant but blind to heterogeneity; CFD sees it but takes days. rCFD built on
CFD can run alongside the process as a soft sensor ($T_{\max}$, $\sigma_T$), aligned with live data
to stop error accumulation. My thesis adds two warnings: a boundary condition can switch the flow
regime, and the heat budget must use the same fluid properties as the recording.
\end{qabox}

\begin{qabox}[Pirker]{\pred CFD basics he may check quickly.}
\from{Numerical Methods in Fluid Mechanics (Pirker) · GovEq, FV, SIMPLE, turbulence · predicted}
\ans \textsf{Re} $=uL/\nu$ (mine 99, laminar). \textsf{Material derivative} $\mathrm{D}/\mathrm{D}t=\partial_t+\mathbf{u}\cdot\nabla$.
\textsf{FVM}: integrate over cells, fluxes through faces, locally conservative, $a_Pc_P=\sum a_{NB}c_{NB}+b$.
\textsf{SIMPLE}: guess $p$, solve $\mathbf{u}^*$, pressure correction from continuity, correct, repeat.
\textsf{Turbulence}: cascade to Kolmogorov $\eta=(\nu^3/\varepsilon)^{1/4}$; RANS models all, LES the
small eddies, DNS none; $\nu_t=C_\mu k^2/\varepsilon$. \textsf{Boussinesq}: constant density except in
the gravity term (mine: 1.5 per cent density change).
\end{qabox}
